% En el preámbulo: \usepackage{booktabs,longtable,array,url}
\begingroup
\small
\renewcommand{\arraystretch}{1.2}
\setlength{\tabcolsep}{4pt}
\begin{longtable}{@{}
  >{\raggedright\arraybackslash}p{0.09\linewidth}
  >{\raggedright\arraybackslash}p{0.29\linewidth}
  >{\raggedright\arraybackslash}p{0.39\linewidth}
  >{\raggedright\arraybackslash}p{0.15\linewidth}@{}}
\caption{Diseño de endpoints principales}
\label{tab:endpoints-principales}\\
\toprule
\textbf{Método} & \textbf{Ruta} & \textbf{Propósito y parámetros} & \textbf{Respuestas} \\
\midrule
\endfirsthead
\multicolumn{4}{c}{\tablename\ \thetable\ (continuación)}\\
\toprule
\textbf{Método} & \textbf{Ruta} & \textbf{Propósito y parámetros} & \textbf{Respuestas} \\
\midrule
\endhead
\midrule
\multicolumn{4}{r}{Continúa en la siguiente página.}\\
\endfoot
\bottomrule
\endlastfoot
POST & \path{/api/v1/login} &
Autentica mediante \texttt{correo} y \texttt{contrasena}. Devuelve un token JWT y los datos del usuario. &
200, 400, 401, 429 \\
GET & \path{/api/v1/pacientes} &
Lista los pacientes autorizados. Parámetros de consulta opcionales: \texttt{page}, \texttt{size} y \texttt{search}. &
200, 401, 403 \\
POST & \path{/api/v1/pacientes} &
Registra un paciente. Recibe \texttt{nombre}, \texttt{correo} y los demás campos permitidos por el catálogo y el rol. &
201, 401, 403, 422 \\
GET & \path{/api/v1/radiographs} &
Consulta solicitudes autorizadas. Filtros opcionales: \texttt{paciente}, \texttt{estado} y \texttt{page}. &
200, 401, 403, 422 \\
POST & \path{/api/v1/radiographs} &
Crea una solicitud con \texttt{paciente\_id}, \texttt{tipo} e \texttt{indicaciones}. El administrador también indica \texttt{odontologo\_id}. &
201, 401, 403, 404, 422 \\
POST & \path{/api/v1/radiographs/:id/access} &
Genera enlace temporal y QR. Recibe el identificador de la solicitud y \texttt{horas}, entre 1 y 72; el valor predeterminado es 24. &
201, 401, 403, 404, 409, 422 \\
GET & \path{/api/v1/radiographs/:id/file} &
Abre el archivo de una solicitud autorizada. El parámetro opcional \texttt{download=1} solicita la descarga. &
200, 302, 401, 403, 404, 422, 503 \\
POST & \path{/api/v1/external/radiographs/access} &
Valida el \texttt{token} del enlace de CERPAX, sin iniciar sesión. Devuelve datos mínimos del estudio y su vencimiento. &
200, 410, 429 \\
POST & \path{/api/v1/external/radiographs/upload-url} &
Prepara la carga mediante URL firmada. Recibe \texttt{token}, \texttt{nombre}, \texttt{tipo} MIME y \texttt{tamano} en bytes. &
200, 410, 422, 429, 503 \\
POST & \path{/api/v1/external/radiographs/upload-complete} &
Confirma la carga en Supabase Storage. Recibe \texttt{token}, \texttt{path}, \texttt{nombre}, \texttt{tipo} y \texttt{tamano}. Completa la solicitud. &
201, 409, 410, 422, 429 \\
\end{longtable}
\noindent\textit{Fuente: Elaboración propia a partir de las rutas implementadas en el sistema.}

\medskip
\noindent Las rutas internas protegidas requieren el encabezado
\texttt{Authorization: Bearer <JWT>} y los permisos correspondientes.
Las rutas externas de CERPAX utilizan un token temporal.
Un origen no autorizado puede generar una respuesta 403; un error interno
puede generar una respuesta 500. El código 400 corresponde, entre otros casos,
a campos de autenticación ausentes o de tipo incorrecto; las validaciones
de datos clínicos utilizan principalmente 422.
\endgroup
